\section{Proofs for Rank-Stratified Geometry}
\label{app:stratified-proofs}

\begin{proof}[Proof of \cref{thm:rank-stratification}]
Let \(\mathrm{Fr}_m\) be the manifold of full-column-rank \(d\times m\)
frames.  Its right \(GL(m)\) action is free and proper, with quotient
\(\Gr(m,d)\).  The diagonal action on
\(\mathrm{Fr}_m\times\SPD^m\) used in
\eqref{eq:total-descriptor-space} is also free.  It is proper by the
sequential criterion: if \((A_k,R_k)\) and
\((A_kB_k,B_k^\top R_kB_k)\) both converge inside the product, then smooth
left inverses of the full-rank \(A_k\) show that \(B_k\) converges; the limit
is invertible because the limiting transformed frame has rank \(m\).
Consequently the quotient is a smooth manifold.  Equivalently, it is the
smooth open cone bundle
\[
  \{(W,g):W\in\Gr(m,d),\ g\text{ is a positive-definite bilinear form on }W\}
\]
over the Grassmannian.  The identification sends \([A,R]\) to
\(W=\range(A)\) and the form defined by
\(g(Ax,Ay)=x^\top Ry\).  This formula is independent of the representative.
The regular descriptor space is the open subbundle on which
\(g-P_0|_W\) is nondegenerate.

Gauge invariance follows because the replacement
\((A,R)\mapsto(AB,B^\top RB)\) leaves
\(P_0AR_0^{-1}(R-R_0)R_0^{-1}A^\top P_0\) unchanged.  Moreover,
\begin{equation}
  \Psi_{P_0,A}(R)-P_0
  =P_0AR_0^{-1}(R-R_0)R_0^{-1}A^\top P_0.
  \label{eq:rank-factorization-proof}
\end{equation}
If \(R-R_0\) is invertible, this difference has rank \(m\).

Conversely, let \(P\in\mathfrak S_m(P_0)\), put \(H=P-P_0\), and define
\(\mathcal W=P_0^{-1}\range(H)\).  Choose full-column-rank \(A\) spanning
\(\mathcal W\).  Then \(P_0A\) spans \(\range(H)\).  Since \(H\) is
symmetric of rank \(m\), there is an invertible symmetric \(C\) such that
\[
  H=P_0ACA^\top P_0.
\]
For example, with \(B=P_0A\), take
\[
  C=(B^\top B)^{-1}B^\top H B(B^\top B)^{-1}.
\]
The identity follows because the range and co-range of \(H\) both equal
\(\range(B)\).  Setting \(R=A^\top PA\) and \(R_0=A^\top P_0A\) gives
\[
  R-R_0=A^\top H A=R_0CR_0,
\]
which is invertible, and \eqref{eq:rank-factorization-proof} reconstructs
\(P\).

The subspace \(\mathcal W\) is intrinsic.  Two frames for it differ by a
unique \(B\in GL(m)\), and their compressed matrices differ by congruence.
This proves bijectivity.  It remains to verify the smooth structures rather
than only the set-theoretic inverse.  After whitening \(P_0\), the bundle
description above identifies a regular descriptor with a pair \((W,T)\),
where \(T\) is a positive-definite self-adjoint endomorphism of \(W\) and
\(T-I_W\) is invertible.  In a local orthonormal frame \(U(W)\), the map is
\[
  (W,T)\longmapsto I+U(W)([T]_{U(W)}-I)U(W)^\top,
\]
so it is smooth and has rank-\(m\) displacement.  Conversely, on the
fixed-rank symmetric displacement manifold, the nonzero spectrum of
\(H=P-P_0\) remains locally separated from zero.  The range projector
\(\Pi_H\) is therefore smooth by spectral functional calculus.  A local
smooth frame of \(\range(\Pi_H)\), together with the restriction of \(P\) to
that range, gives exactly the local inverse \((W,T)\).  Transition between
two such frames is smooth congruence, which is the quotient transition
already built into the form bundle.  Fixed inertia classes of \(H\) are
open-and-closed components of the fixed-rank symmetric manifold; the same
argument applies on every component, so no signature branch is omitted.
Hence \(\Gamma_{m,P_0}\) and its inverse are smooth, and the map is a
diffeomorphism.

Every SPD matrix lies in exactly one rank stratum.  A limit of rank-\(m\)
symmetric differences has rank at most \(m\).  Conversely, if
\(\rank(P-P_0)=r<m\), whiten by \(P_0\), choose \(m-r\) orthonormal kernel
directions of the whitened difference, and add arbitrarily small nonzero
eigenvalues in those directions.  The perturbation stays SPD for sufficiently
small magnitudes, has rank \(m\), and converges to \(P\).  This proves
\eqref{eq:strata-closure}.
\end{proof}

\begin{proof}[Proof of \cref{thm:stratified-line-element}]
Under \(A^\top P_0A=I\), put \(U=P_0^{1/2}A\).  Whitening sends the base
point and the differential represented by \(\xi_i=(H_i,K_i)\) to
\[
  \Gamma\longleftrightarrow
  \begin{pmatrix}R&0\\0&I\end{pmatrix},
  \qquad
  D\widetilde\Gamma[\xi_i]\longleftrightarrow
  \begin{pmatrix}
    H_i&(R-I)K_i^\top\\
    K_i(R-I)&0
  \end{pmatrix}.
\]
Using
\(
\langle X,Y\rangle_{P,\AI}=\tr(P^{-1}XP^{-1}Y)
\), direct block multiplication for directions \(i,j\) gives
\[
  \tr(R^{-1}H_iR^{-1}H_j)
  +2\tr(K_i(R-I)R^{-1}(R-I)K_j^\top),
\]
which is \eqref{eq:completion-pair-kernel}.  Setting \(i=j\) proves
\eqref{eq:stratified-line-element}.  Factoring its two blocks gives
\eqref{eq:completion-feature}.  The pair matrix is therefore a Gram matrix,
which proves its positive semidefiniteness, the rank identity
\eqref{eq:completion-gram-rank}, and the orthogonality criterion
\eqref{eq:completion-pair-orthogonality}.

The differential vanishes exactly when \(H=0\) and \(K(R-I)=0\), proving
\eqref{eq:stratified-kernel}.  If \(m<d\) and \(R-I\) is singular, a nonzero
row in its nullspace can be placed in \(K\), producing a nonzero kernel
direction; if \(R-I\) is invertible, the kernel condition gives \(K=0\).
Thus for \(m<d\) the pullback is positive definite exactly on the regular
domain of \cref{thm:rank-stratification}.  If \(m=d\), the horizontal
Grassmann block has dimension zero, so \(K\) is absent and the visible term is
positive definite for every \(R\in\SPD^d\).  In this case the base
Grassmannian is a point and, after whitening, the total descriptor is simply a
positive form on the full space; the completion map returns that same form.
It is therefore a global diffeomorphism onto \(\SPD^d\).

The eigenvalues of \(\Xiop(R)\) are
\((\lambda_i(R)-1)^2/\lambda_i(R)\).  Under the spectral assumptions,
\[
  \lambda_{\min}(\Xiop(R))\ge\frac{\zeta^2}{b},
  \qquad
  \lambda_{\max}(\Xiop(R))
  \le\max_{x\in[a,b]}\frac{(x-1)^2}{x}.
\]
Substitution into the exact line element proves
\eqref{eq:stratified-conditioning}.
\end{proof}

\begin{proof}[Proof of \cref{cor:reduced-completion-value}]
The shortest-completion identity \eqref{eq:spd-shortest-completion} gives
\[
  \min_{A^\top PA=R}\frac12\dai(P_0,P)^2
  =\frac12\dai(A^\top P_0A,R)^2.
\]
Under the normalization \(A^\top P_0A=I\), this is
\eqref{eq:reduced-completion-energy} and is independent of the horizontal
subspace coordinate.  Derivatives with respect to that coordinate at fixed
\(R\) consequently vanish.  The minimizer is the explicit section
\(P=\Psi_{P_0,A}(R)\), so evaluating it requires no further vertical
stationarity equation.  Differentiating that section instead measures the
motion of the minimizing ambient matrix, whose AIRM inner product is
\eqref{eq:completion-pair-kernel}; it is not the Hessian of the reduced value.
\end{proof}

\paragraph{Frozen--relaxed decomposition in the two-dimensional example.}
The cancellation mentioned in \cref{rem:value-versus-motion} can be made
explicit in fixed ambient coordinates.  Write
\[
  u(\theta)=(\cos\theta,\sin\theta)^\top,
  \qquad
  P=\begin{pmatrix}a&c\\c&b\end{pmatrix},
\]
and eliminate the constraint \(u(\theta)^\top Pu(\theta)=r\) as
\[
  a=
  \frac{r-2c\sin\theta\cos\theta-b\sin^2\theta}
       {\cos^2\theta}.
\]
For
\(F(b,c,\theta)=\tfrac12\fro{\log P}^2\), the hidden minimizer at
\(\theta=0\) is \((b,c)=(1,0)\).  Direct differentiation there gives
\[
  F_{\theta\theta}=\frac{2(r-1)\log r}{r},
  \qquad
  F_{\theta c}=-\frac{2\log r}{r},
  \qquad
  F_{cc}=\frac{2\log r}{r(r-1)},
\]
with \(F_{\theta b}=F_{bc}=0\) and \(F_{bb}=1\).  The hidden Hessian is
positive definite because \(\log r/(r-1)>0\).  Consequently the relaxed
second derivative is
\[
  F_{\theta\theta}
  -F_{\theta z}F_{zz}^{-1}F_{z\theta}
  =\frac{2(r-1)\log r}{r}
  -\frac{2(r-1)\log r}{r}
  =0,
  \qquad z=(b,c).
\]
Thus the frozen direct term and the nonzero Schur term cancel, consistently
with the constant reduced value.  The response
\(c'(0)=-F_{cc}^{-1}F_{c\theta}=r-1\) is the off-diagonal derivative of the
canonical completion in \cref{rem:value-versus-motion}.

\begin{proof}[Proof of \cref{prop:moving-reference}]
By definition, \(P_t=G_tQ_tG_t^\top\).  Differentiating and multiplying by
\(G_t^{-1}\) and \(G_t^{-\top}\) gives
\[
  G_t^{-1}\dot P_tG_t^{-\top}
  =\dot Q_t+\Omega_tQ_t+Q_t\Omega_t^\top
  =\mathfrak D_tQ_t.
\]
Congruence invariance of AIRM yields \eqref{eq:moving-reference}.

If \(G_t'=G_tO_t\), then
\[
  Q_t'=O_t^\top Q_tO_t,
  \qquad
  \Omega_t'=O_t^\top\Omega_tO_t+O_t^\top\dot O_t.
\]
Differentiating \(Q_t'\) and substituting these identities cancels all terms
containing \(O_t^\top\dot O_t\), leaving
\[
  \mathfrak D_t'Q_t'=O_t^\top(\mathfrak D_tQ_t)O_t.
\]
Orthogonal congruence preserves the displayed Frobenius norm.
\end{proof}

\begin{proof}[Proof of \cref{cor:rank-loss}]
After whitening, the nonzero singular values of
\[
  P_t-I=U_t(R_t-I)U_t^\top
\]
are exactly those of \(R_t-I\).  Singular values vary continuously with the
matrix.  If the limit has rank at most \(m-r\), at least \(r\) of the first
\(m\) singular values converge to zero.  The same therefore holds for
\(R_t-I\).
\end{proof}
